% AUTHOR: Amlal El Mahrouss
% PURPOSE: A Short Reference for Algebraic Properties.

\documentclass[10pt]{article}

%% ---- Core Packages ----------------------------------------
\usepackage[T1]{fontenc}
\usepackage[utf8]{inputenc}
\usepackage{lmodern}
\usepackage{microtype}

%% ---- Page Geometry ----------------------------------------
\usepackage[
top=1in, bottom=1in,
left=0.75in, right=0.75in,
columnsep=0.25in
]{geometry}

%% ---- Mathematics ------------------------------------------
\usepackage{amsmath, amssymb, amsthm, amsfonts}
\usepackage{mathtools}
\usepackage{bm}
\usepackage{textcomp}
\usepackage{esint}

%% ---- Figures & Tables -------------------------------------
\usepackage{graphicx}
\usepackage{booktabs}
\usepackage{multirow}
\usepackage{array}
\usepackage{caption}
\usepackage{subcaption}
\usepackage{float}

%% ---- Code Listings ----------------------------------------
\usepackage{listings}
\usepackage{xcolor}
\usepackage{algorithm}
\usepackage{algpseudocode}

\lstdefinestyle{codestyle}{
	backgroundcolor=\color{gray!8},
	basicstyle=\ttfamily\footnotesize,
	breakatwhitespace=false,
	breaklines=true,
	captionpos=b,
	commentstyle=\color{green!50!black},
	keywordstyle=\color{blue}\bfseries,
	stringstyle=\color{orange!70!black},
	numberstyle=\tiny\color{gray},
	numbers=left,
	numbersep=5pt,
	showstringspaces=false,
	frame=single,
	rulecolor=\color{gray!40},
	tabsize=2
}
\lstset{style=codestyle}

%% ---- Hyperlinks -------------------------------------------
\usepackage[
colorlinks=true,
linkcolor=blue!70!black,
citecolor=green!50!black,
urlcolor=blue!60!black
]{hyperref}
\usepackage{url}

%% ---- Bibliography -----------------------------------------
\usepackage[numbers, sort&compress]{natbib}

%% ---- Miscellaneous ----------------------------------------
\usepackage{enumitem}
\usepackage{siunitx}
\usepackage{cleveref}

%% ---- Theorem Environments ---------------------------------
\theoremstyle{plain}
\newtheorem{theorem}{Theorem}[section]
\newtheorem{lemma}[theorem]{Lemma}
\newtheorem{corollary}[theorem]{Corollary}
\newtheorem{proposition}[theorem]{Proposition}

\theoremstyle{definition}
\newtheorem{definition}[theorem]{Definition}
\newtheorem{example}[theorem]{Example}

\theoremstyle{remark}
\newtheorem{remark}[theorem]{Remark}

%% ---- Custom Commands --------------------------------------
\newcommand{\R}{\mathbb{R}}
\newcommand{\N}{\mathbb{N}}
\newcommand{\Z}{\mathbb{Z}}
\newcommand{\E}{\mathbb{E}}
\newcommand{\Prob}{\mathbb{P}}
\newcommand{\norm}[1]{\left\lVert #1 \right\rVert}
\newcommand{\abs}[1]{\left\lvert #1 \right\rvert}
\newcommand{\inner}[2]{\left\langle #1,\, #2 \right\rangle}
\newcommand{\ie}{\textit{i.e.}\xspace}
\newcommand{\eg}{\textit{e.g.}\xspace}
\newcommand{\etal}{\textit{et al.}\xspace}

%% ---- Caption Formatting -----------------------------------
\captionsetup{
	font=small,
	labelfont=bf,
	format=hang,
	justification=justified
}

\renewcommand{\qedsymbol}{\ensuremath{\blacksquare}}

\usepackage{lettrine}

%% ---- Header / Footer --------------------------------------
\usepackage{fancyhdr}
\pagestyle{fancy}
\fancyhf{}
\fancyhead[R]{\small\thepage}
\renewcommand{\headrulewidth}{0.4pt}

%% ============================================================
%%  FRONT MATTER
%% ============================================================

\title{%
	\vspace{-1.5em}%
	\large\textbf{Little Learner Reference Formulas.}%
}
\author{By A. E. Mahrouss}
\date{\today}

%% ============================================================
\begin{document}
%% ============================================================

\maketitle
\tableofcontents
\newpage

\section{Synopsis:}

\lettrine[lines=1, lhang=0.1, loversize=0.2]{L}ET the following notes on Little Learner's state the tactic on its training. LiL adopts a dynamic approach to training. Letting it be able to unlearn and relearn missing or wrong concepts.

\section{Introduction and Definitions:}

One may ask for definitions and rigorous mathematical grounding before proceeding development of LiL. The two sub-sections provides exactly that.\\
One may note that the remarks and theorems are currently being benchmarked, addendum will be added accordingly.

\subsection{Definitions:}

One may train LiL using the following theorems and remarks:

\begin{theorem}
Let the following variables in $\mathbb{R}$ be defined as follows:
\begin{align*}
    \mathcal{X}_n &= \sum_{k=1}^{+\infty} \mathcal{T}_k \otimes \mathcal{Y}_k \otimes C_k.
\end{align*}
And for the training nabla operator in $\mathbb{R}$ given for a single coordinate system of $x$.:
\begin{align*}
	\nabla \mathcal{X}_n &= i \cdot \frac{\partial x}{\partial \, (\sum_{k=1}^{+\infty} \mathcal{T}_k \otimes \mathcal{Y}_k \otimes C_k)} + F(x).
\end{align*}
Where $C_k$ is a non-zero constant term in $\mathbb{R}$. Where $\mathcal{X}$ is a sum of tensor products of $\mathcal{T}_k$ and $\mathcal{Y}_k$ for a sequence $k=1, 2, \ldots$ that holds for $\mathbb{N}$.\\
Such that $\forall F(x) \in \mathbb{R}$ and $F(x) \subset \mathcal{T}$.
\end{theorem}

\begin{remark}
    $\mathcal{L}$ shall be non-zero and non-negative, as it is a sum of positive terms in $\mathbb{R}$.\\
	Inserting the theorem to LiL' training tree alongside an hash-value pair. One may be able to query for the theorem.
\end{remark}

\begin{remark}
	Another example of a definition this time, would be a Taylor Series for $f(x) \in \mathbb{R}$ and $a \in \mathbb{R}$, which states:
	\begin{align*}
		\mathcal{T}(x) &= \sum_{n=0}^{\infty} \frac{f^{(n)}(a)}{n!} (x-a)^n.
	\end{align*}
	For $\mathcal{T}(x)$ representing the Taylor Series expansion of $f(x)$ about the point $a$.
\end{remark}

\section{Going Forward:}

\lettrine[lines=1, lhang=0.1, loversize=0.2]{N}EXT we will observe and study several tactics we can consider for LiL's Input/Output training.

\nocite{*}

\bibliographystyle{acm}
%% \bibliography{}

%% ============================================================
\end{document}
%% ============================================================
